\documentclass[../DoAn.tex]{subfiles}
\begin{document}

\begin{center}
  \Large{\textbf{TÓM TẮT}}\\
\end{center}
\vspace{1cm}
Trong bối cảnh các vấn đề xã hội ngày càng phức tạp, tình trạng người thân bị mất tích do nhiều nguyên nhân như tai nạn, bệnh lý hoặc các yếu tố bất khả kháng đã trở thành một vấn nạn nhức nhối. Các phương pháp tìm kiếm truyền thống như trình báo công an, đăng tin trên các phương tiện truyền thông thường tốn nhiều thời gian, thủ tục phức tạp và tiềm ẩn rủi ro bị các đối tượng xấu lợi dụng thông tin để lừa đảo, gây thêm tổn thất cho gia đình nạn nhân.

Đồ án này đề xuất một khung kiến trúc toàn diện nhằm giải quyết vấn đề trên, xây dựng Hệ thống phát hiện người mất tích dựa trên dữ liệu ảnh chụp – một mô hình tập trung, an toàn và hiệu quả, áp dụng công nghệ Trí tuệ Nhân tạo (AI) để kết nối gia đình và người tìm thấy.

Cốt lõi của hệ thống là việc áp dụng thuật toán Deep Learning cho công nghệ nhận dạng khuôn mặt, và đạt độ chính xác ổn định, kết hợp với một kiến trúc nghiệp vụ an toàn. Điểm độc đáo của mô hình không nằm ở bản thân thuật toán nhận dạng, mà ở quy trình xử lý khép kín và bảo mật.

An ninh và xác thực là ưu tiên hàng đầu của hệ thống nhằm ngăn chặn tuyệt đối các hành vi xâm nhập và lạm dụng. Quy trình được thiết kế để bảo vệ tất cả các bên liên quan: Tiếp nhận thông tin, Xác thực thông tin, Xử lý và Thông báo.

Mô hình được đề xuất không chỉ giải quyết được cả bốn trường hợp nghiệp vụ chính (gia đình đăng tin trước, người tìm thấy đăng tin trước, đăng tin trùng lặp, và kẻ xấu mạo danh) mà còn mang lại sự thuận tiện, giảm thiểu thủ tục hành chính cho người dân.
\begin{flushright}
  % Sinh viên thực hiện\\
  \begin{tabular}{@{}c@{}}
    % \textit{(Ký và ghi rõ họ tên)}
  \end{tabular}
\end{flushright}

\end{document}
